\documentclass[10pt,a4paper]{amsart}
\usepackage[margin=19mm]{geometry}
\usepackage{amsmath,amssymb,mathtools}
\usepackage[scr=rsfs,cal=euler]{mathalfa}
\usepackage{xfrac}
\usepackage[unicode,hidelinks,bookmarks=false]{hyperref}
\newcommand{\ie}{i.e.\ }
\newcommand{\cf}{cf.\ }
\newcommand{\resp}{respectively\ }
\newcommand{\Subsection}{\subsection}
\providecommand{\nobreakdash}{\nobreak\hbox{-}}
\setlength{\emergencystretch}{2em}
\multlinegap0pt
\usepackage{mathrsfs}
\newtheorem{theorem}{Theorem}[section]
\newtheorem{proposition}[theorem]{Proposition}
\newtheorem{lemma}[theorem]{Lemma}
\newtheorem{corollary}[theorem]{Corollary}
\newtheorem{remark}[theorem]{Remark}
\newcommand\R{\mathbb{R}}
\begin{document}
\title[Propagation of Regularity for Schroedinger Equations with Ti]{Propagation of Regularity for Schroedinger Equations with Time Dependent Potentials}
\author{Avy Soffer}
\date{}
\maketitle
\noindent\emph{Source and licence.} Propagation of Regularity for Schroedinger Equations with Time Dependent Potentials. Version arXiv:2605.27321v3 (CC BY 4.0). This material is distributed under the Creative Commons Attribution 4.0 International licence, http://creativecommons.org/licenses/by/4.0/, as recorded in the arXiv OAI licence metadata for this exact version and shown on the abstract page https://arxiv.org/abs/2605.27321v3. Full rights notice: licenses/arxiv-2605.27321-CC-BY-4.0.txt.\par\smallskip
\noindent\textit{Editorial scope.} Counted unit: the original \texttt{theorem} environment \texttt{-} at source lines 1200--1219 together with its complete proof at lines 1221--1362; the delivered document keeps source lines 1198--1363 so that every referenced label and every citation resolves inside it. Native editable source is the paper's own concatenated TeX; the AMS wrapper is editorial formatting only. No publisher class, upstream script or unrelated artwork is redistributed.
\section{$H^{2}$ - Regularity on the Outgoing Waves}


\begin{theorem}

Let $\psi(t)$ be a solution of the SE as considered above.
Assume the initial condition is well localized in space and in $H^{2}\left(R^{3}\right)$.


Then,\newline
a) If $\sup _{t}\langle\phi(t), H(t) \phi(t)\rangle \lesssim O(1)$
where $\phi(t)=F\left(\frac{A}{t^{\alpha}} \geqslant 1\right) \psi(t)$,
then

\begin{equation}
\int_{1}^{T}\left\langle\Delta \psi(t),\left[\left\langle\frac{A}{t^{\alpha}}\right\rangle \ln \left\langle A / t^{\alpha}\right\rangle\right]^{-1} \Delta \psi(t)\right\rangle \frac{d t}{t^{\alpha}} \leq O(1) \tag{C.2}
\end{equation}

b) If for suitably localized $G(A)$ \begin{equation}
\sup_t \langle \nabla\phi(t),G(A)\langle\frac{A}{t^{\alpha}}\rangle\ln{\langle\frac{A}{t^{\alpha}}\rangle}  \nabla \phi(t)\rangle\leq O(1) \tag{C.3a}\end{equation}
then \begin{equation}
\int_{1}^{T} \frac{d t}{t}\left\langle\Delta \psi(t), G(A) F\left(\frac{A}{t^{\alpha}} \geqslant 1\right) \Delta \psi(t)\right\rangle \leq O(1). \tag{C.3b}\end{equation}  
\end{theorem}

\begin{proof}

\begin{align*}
\partial_{t}\langle\psi(t), & \left.F_{A, M}\left(\frac{A}{M t^{\alpha}} \geq 1\right) H(t) F_{A, M}\left(\frac{A}{M t^{\alpha}} \geq 1\right) \psi(t)\right\rangle \\
= & \left\langle\frac{1}{M t^{\alpha}} \tilde{F}_{A, M}^{\prime}\left[2 p^{2}-\frac{\alpha A}{t}\right] H(t) F_{A, M}+c . c .\right\rangle \\
& +\left\langle\psi(t), F_{A, M} \frac{\partial V}{\partial t} F_{A, M} \psi(t)\right\rangle  \tag{C.4}\\
& +\left\langle\psi(t),\left(i\left[V, F_{A, M}\right] H(t) F_{A, M}+c . c .\right) \psi(t)\right\rangle \equiv I_{1}+I_{2}+I_{3}
\end{align*}


We proved in the previous section estimates away from the PS.

So, to deal with the first term, we write

\begin{align*}
& {\left[2 p^{2}-\frac{\alpha A}{t}\right]=\left[2 p^{2}-\frac{\alpha A}{t}\right] F_{A, C}^l\left(\frac{2 p^{2} t-A}{t^{\alpha}} \geqslant 1\right) } \\
+ & {\left[2 p^{2}-\frac{\alpha A}{t}\right] F_{A, C}^L\left(\frac{A-2 p^{2} t}{t^{\alpha}} \geqslant 1\right) } \\
+ & {\left[2 p^{2}-\frac{\alpha A}{t}\right] F_{A, C}^0\left(\frac{\left|A-2 p^{2} t\right|}{t^{\alpha}} \leq 1\right) } \tag{C.5}
\end{align*}


The first and third terms in equation  C.5 are positive (for $2p^2 t\geq t^{\alpha}$).
The second term is supported where $\frac{A}{t} \geqslant 2 p^{2}+t^{\alpha-1}$ and multiplied by $\tilde{F}_{A, M}^{'}$ where $\frac{A}{t^{\alpha}} \sim M$.

Therefore, $M t^{\alpha-1} \geqslant 2 p^{2}+t^{\alpha-1}$.

Consequently,

\begin{align*}
& \left\langle\tilde{F}_{A, M}^{\prime}\left[2 p^{2}-\frac{\alpha A}{t}\right] F_{A, C}^L\left(\frac{A-2 p^{2} t}{t^{\alpha}} \geqslant 1\right) H(t) F_{A, M}\right) \\
\lesssim & \left\langle\tilde{F}_{A, M}^{\prime} M^{2} t^{2 \alpha-2} F_{A, C}^L F_{A, M}\right\rangle  \tag{C.6}\\
& \text { since } F_{2}\left(p^{2} \leq M t^{\alpha-1}\right) H(t) \lesssim F_{2} M t^{\alpha-1}+F_{1} V
\end{align*}

and
$$
\begin{aligned}
\tilde{F}_{A, M}^{\prime} \cdot F_{1} V \sim & \tilde{F}_{A, M}^{\prime}\langle A\rangle^{-2} F_{1} A^{2} V \sim \\
& \left\langle M t^{\alpha}\right\rangle^{-2} M t^{\alpha-1}\left( |x|^{2} V\right) \sim M^{-1} t^{-1-\alpha}.
\end{aligned}
$$

So,

\begin{align*}
I_{1} &\geqslant\langle (Mt^{\alpha})^{-1}\tilde F'_{A,M}2p^4 F_{A,C}^lF_{A,M}\rangle-\left\langle\left(M t^{\alpha}\right)^{-1} \tilde F_{A, M}^{'} M^{2} t^{2 \alpha-2} F_{A, C}^l\left(\frac{A-2 p^{2} t}{t^{\alpha}} \geqslant 1\right) F_{A, M}\right\rangle \\
I_{2} & =\left\langle\psi(t), F_{A, M} \frac{\partial V}{\partial t} F_{A, M} \psi(t)\right\rangle  \tag{C.7}\\
& \approx\left\langle\psi(t), F_{A, M}\langle A\rangle^{-2}\left(A^{2}+1\right)\langle x\rangle^{-m}\left(A^{2}+1\right)\langle A\rangle^{-2} F_{A, M} \psi(t)\right\rangle  \tag{C.8}\\
& \approx\left\langle\Delta \psi(t),\langle x\rangle^{-\sigma} \tilde{F}_{A,M} t^{-4 \alpha}\langle x\rangle^{-m+4+2 \sigma} \tilde{F}_{A, M}\langle x\rangle^{-\sigma} \Delta \psi(t)\right\rangle M^{-4} \\
& \approx O\left(t^{-4 \alpha}\right)\|\Delta \psi(0)\|_{2}^{2} M^{-4},
\end{align*}

provided $m \geqslant 4+2 \sigma ; \quad \psi(0) \in H^{2}\left(\mathbb{R}^{3}\right)$.
\begin{remark}
The above estimate with $m\geqslant 4+2 \sigma$ is not necessary, and $\geqslant 2+2 \sigma$ is sufficient.
In the above we demonstrate the maximal decay possible at the $H^2$ level.
\end{remark}

The control of  $I_3$ is similar:
In this case, due to the $H(t)$ factor, we can only multiply by $\langle A\rangle^{2}$ and divide.
By restricting the domain to the PS or where
$A \geqslant 2 p^{2} t$, two powers of $\langle A\rangle ^{-1}$ give a decay of order $t^{-2}$, so, such regions are bounded by
$$
\|\psi(0)\|_{H^{2}}^{2} t^{-2}\left(M t^{2}\right)^{-1} .
$$

In the region $A \leq 2 p^{2} t$, the $\langle A\rangle^{-2}$ gives
$$
\|\psi(0)\|_{H^{2}}^{2}\left(M t^{\alpha}\right)^{-2}
$$

We also used the Commutator Expansion Lemma for
$$
\left[V, F_{A, M}\right] \sim x \cdot \nabla V F_{A, M}^{\prime}\left(M t^{\alpha}\right)^{-1}+\text { h.o.t. }
$$

So,

\begin{equation*}
I_{3} \lesssim\left[\left(M t^{\alpha}\right)^{-3}+t^{-2}\left(M t^{\alpha}\right)^{-1}\right]\|\psi(0)\|_{H^{2}}^{2} \tag{C.9}
\end{equation*}


Combining the equations C.4-C.9
we get

\begin{align*}
& \left\langle F_{A, M} H(t) F_{A, M}\right\rangle_{t=T}-\left\langle F_{A, M} H(t) F_{A, M}\right\rangle_{t=1} \\
& =\int_{1}^{T} \frac{d t}{t^{\alpha}}\left\langle\frac{1}{M} \tilde{F}_{A, M}^{\prime} 2 p^{2} F_{A, C}^{+} p^{2} \tilde{F}_{A, M}^{\prime}\right\rangle+h.o.t. \\
& +\int_{1}^{T} t^{2\alpha} \frac{d t}{t^{2} t^{\alpha}} \frac{M^{2}}{M}\left\langle\tilde{F}_{A, M}^{\prime} F_{2}\left(p^{2} \leq M t^{\alpha-1}\right) F_{A, C}^{-} F_{A, M}+c . c .\right\rangle \\
& +\int_{1}^{T} d t\left\{O\left(t^{-4 \alpha}M^{-4}\right)+\left(M t^{\alpha}\right)^{-3}+O\left(t^{-2}\right)\left(M t^{\alpha}\right)^{-1}\right\}^{2} \|\Delta \psi(0)\|_{2}^{2}. \tag{C.10}
\end{align*}


Here,
$$
F_{A, C}^{+}=F_{A, C}\left(\frac{\left|2 p^{2} t-A\right|}{t^{\alpha}} \leq 1\right)+F_{A, C}\left(\frac{2 p^{2} t-A}{t^{\alpha}} \geqslant 1\right)
$$

We use $\tilde{F}_{A, M}^{\prime}$ to denote generic function of the form $\widetilde{F}_{A, M}\left(\frac{A}{M t^{\alpha}}=1\right)$.

Next, we will use that
$$
\sum_{n=0}^{\infty}\left(M_{0} 2^{n}\right)^{l} F_{A, M_{0} 2^{n}} \sim\left(\frac{A}{t^{\alpha}}\right)^{l} \ln \left(\langle A\rangle / t^{\alpha}\right) F_{A}\left(\frac{A}{t^{\alpha}} \geqslant M_{0}\right)
$$
and
$$
\sum_{n=0}^{\infty}\left( M_02^{n}\right)^{l} / n F_{A, M_0 2^{n}} \sim\left(\frac{A}{t^{\alpha}}\right)^{l} F_{A}\left(\frac{A}{t^{\alpha}} \geq M_0\right)
$$



We sum over $M=M_{0} 2^{n}, n=0,1,2, \ldots$ both sides of eq. C.10, and sum over $n$ up to $M\sim cT^{1-\alpha}.$

{\bf The LHS becomes}
$$
\left\langle  F_{A}\left(cT^{1-\alpha}\geq\frac{A}{t^{\alpha}} \geq M_{0}\right) H(t) F_{A}\left(cT^{1-\alpha}\geq\frac{A}{t^{\alpha}} \geq M_{0}\right) \ln T\right\rangle_{t}-\left\langle  F_{A}\left(cT\geq A \geqslant M_{0}\right) \ln T \, H(1) F_{A}\left(cT\geq A\geq M_{0}\right)\right\rangle.
$$

{\bf The RHS}

The first term
$$
 \int_{1}^{T} d t\left\langle \langle A\rangle^{-1} F_{A}\left(cT^{1-\alpha}\geq \frac{A}{t^{\alpha}} \geqslant M_{0}\right) p^{2} F_{A,C}^{+} p^{2} F_{A}\left(cT^{1-\alpha}\geq\frac{A}{t^{\alpha}} \geqslant M_{0}\right)\right\rangle\sim
 $$
$$
\frac{1}{cT}\int_1^T \left \langle F_A p^2 F^+_{A,C} p^2 F_A\right \rangle d t.
$$


The second term
$$
\int_{1}^{T} \frac{d t}{t^{2}}\left\langle A F_{A}\left(cT^{1-\alpha}\geq\frac{A}{t^{\alpha}} \geqslant M_{0}\right) F_2 F_{A ,C}^{-} F_{A}\left(cT^{1-\alpha}\geq \frac{A}{t^{\alpha}} \geqslant 1\right)\right\rangle
$$

The third term
$$
C+\int_{1}^{T} \frac{d t}{t^{2+\alpha}}\|\psi(0)\|_{H^{2}}^{2}
$$

The LHS is uniformly bounded by $\ln T$ since $A\leq cT.$
We conclude that a sequence of times going to infinity exists such that the $H^2$ norm remains uniformly bounded on the phase space support of the localization of $A.$ This bound is then extended to all times, using as PROB $p^2G(A,t)p^2$ with $G$ localizing at the relevant domain in phase space. \end{proof}


\end{document}
